\textbf{Segmentation backbones and losses.} U-Net~\cite{ronneberger2015u} was designed for segmentation from few annotated images with heavy augmentation, which makes it the natural model for a data-size study; we use a much smaller variant. Overlap-based losses such as the Dice loss of V-Net~\cite{milletari2016v} are common for class imbalance; we keep pixel-wise losses so that the noise-robust variants are directly comparable, and leave Dice-type losses out of scope.

\textbf{Memorisation and noisy labels.} Deep networks can fit random labels~\cite{zhang2016understanding}, but on real data they tend to learn simple, clean patterns before memorising noise~\cite{arpit2017closer}; classification accuracy can remain high under massive uniform noise if the dataset is large enough~\cite{rolnick2017deep}. This motivates our interaction between noise and data size and our training-length sweep. Surveys of the field are given by Song et al.~\cite{song2020learning}.

\textbf{Noise-robust training.} GCE~\cite{zhang2018generalized} interpolates between CE and the mean absolute error, SCE~\cite{wang2019symmetric} adds a reverse cross-entropy term, Co-teaching~\cite{han2018co} trains two networks that select small-loss samples for each other, and early-learning regularisation~\cite{liu2020early} exploits the early-learning phenomenon. We reimplement GCE and SCE (they are one-line pixel-wise losses and need no extra network); Co-teaching and ELR need sample-level selection or per-sample memory that does not map cleanly to pixel-wise noise in a small study and are not run.

\textbf{Segmentation under noisy annotation and data scaling.} Karimi et al.~\cite{karimi2019deep} survey techniques for learning with noisy labels in medical image analysis, including segmentation, which is the closest prior work to ours; our contribution is the controlled synthetic setting where noise type, level and training-set size are crossed with exactly known ground truth. Empirical data-scaling studies~\cite{hestness2017deep} motivate reading performance as a function of $N$.
